% J25 — The CL_TSML Composition Lattice on Z/10Z:
%        Structural Axioms, Independence, and a 73-HARMONY Forcing Theorem
%
% Brayden R. Sanders, M. Gish (Sanders + Gish lane, per AUTHOR_LANES_v2)
% Submitted to: Algebraic Combinatorics
% Source corpus: Atlas/LENS_TAXONOMY_2026-05-06/CL_FORCING_AXIOMS.md (2026-05-06)
% MSC 2020: 08A40, 20N02, 20N05
% Tier: B (axiomatization of an explicit finite commutative magma)
% Closest published precedent: Drápal-Wanless (2021), JCT-A 184, 105510

\documentclass[11pt,reqno]{amsart}

\usepackage{amsmath, amssymb, amsthm, mathtools}
\usepackage[margin=1in]{geometry}
\usepackage[colorlinks=true,linkcolor=blue,citecolor=blue,urlcolor=blue]{hyperref}
\usepackage{microtype}
\usepackage{booktabs}
\usepackage{array}

\theoremstyle{plain}
\newtheorem{theorem}{Theorem}[section]
\newtheorem{proposition}[theorem]{Proposition}
\newtheorem{lemma}[theorem]{Lemma}
\newtheorem{corollary}[theorem]{Corollary}

\theoremstyle{definition}
\newtheorem{definition}[theorem]{Definition}
\newtheorem{conjecture}[theorem]{Conjecture}

\theoremstyle{remark}
\newtheorem*{remark}{Remark}

\newcommand{\Z}{\mathbb{Z}}
\newcommand{\BUMP}{\mathrm{BUMP}}
\newcommand{\TSML}{\mathrm{TSML}}
\newcommand{\BHML}{\mathrm{BHML}}
\newcommand{\CL}{\mathrm{CL}}

\title[The $\CL_{\TSML}$ composition lattice]
{The $\CL_{\TSML}$ Composition Lattice on $\Z/10\Z$: \\
 Structural Axioms, Independence, and a $73$-HARMONY Forcing Theorem}

\author{Brayden R.\ Sanders}
\address{7Site LLC, Hot Springs, Arkansas, USA}
\email{brayden@7site.co}

\author{M.\ Gish}
\address{Independent Researcher, Hot Springs, Arkansas, USA}
\email{monica.gish1992@gmail.com}

\subjclass[2020]{08A40, 20N02, 20N05}
\keywords{commutative magma, composition lattice, absorbing element,
near-absorbing element with puncture, exceptional cells, axiom
independence, finite non-associative algebra, $\Z/10\Z$ substrate}

\date{\today}

\begin{document}

\begin{abstract}
We isolate a list of seven \emph{structural} axioms~$S_{1}$--$S_{7}$
on a commutative magma $M$ on the carrier $\Omega = \Z/10\Z$, each
axiom expressing a property of $M$ rather than a list of cell values.
Together with the explicit specification of two distinguished
elements ($V = 0$ playing the role of a near-absorber with a single
puncture, and $H = 7$ playing the role of an absorber) and an
explicit list of five \emph{exceptional positions} with their values
($S_{7}$), the seven axioms uniquely determine a $10\times 10$
multiplication table $\CL_{\TSML}$ on $\Omega$ with $73$ HARMONY
($= 7$) cells, $17$ VOID ($= 0$) cells, and $10$ exceptional cells
distributed over the $5$ enumerated positions. We prove the forcing
theorem (Theorem~\ref{thm:forcing}) by direct cell-counting, and we
prove the independence of $S_{1}$--$S_{7}$ in the sense that for each
axiom $S_{i}$ there exists a commutative magma satisfying
$\{S_{j} : j \neq i\}$ but not $S_{i}$ (Theorem~\ref{thm:independence}).
The $\CL_{\TSML}$ lattice sits in the same intellectual neighborhood
as the maximally-non-associative quasigroups of
\cite{DrapalWanless2021}: a small finite commutative non-associative
structure with explicit invariants, in our case characterized by
near-absorption with puncture rather than maximal associativity
defect. The axiom system places $\CL_{\TSML}$ on the same axiomatic
footing as classical absorbing-element semigroups
\cite{HowieSemigroup, Clifford1961} while making explicit the small
finite enumeration of exceptional positions needed to specify a
non-trivial commutative magma on the punctured-VOID,
HARMONY-saturated $\Z/10\Z$ substrate.

\smallskip\noindent
\textbf{Scope and tier discipline.} \emph{Proven.} Axioms
$S_{1}$--$S_{7}$ uniquely determine $\CL_{\TSML}$ as a $10 \times 10$
table (Theorem~\ref{thm:forcing}); the axiom system is non-redundant
in the sense that for each $i$, an explicit witness magma
$M_{i}$ violates $S_{i}$ (Theorem~\ref{thm:independence}); $M_{2}$
and $M_{7}$ are strict independence witnesses (failing only the
targeted axiom), while the other five witnesses fail $S_{i}$
together with one or two additional axioms due to the tight
combinatorial interaction of the count constraint $S_{5}$, the
position constraint $S_{6}$, and the puncture constraint $S_{2}$.
\emph{Computed.} A reference Python script (\texttt{cl\_forcing.py},
runtime $< 1$\,s) verifies the $100$-cell match, computes the
HARMONY/VOID/exceptional partition, and instantiates seven witness
magmas with their full axiom-status vectors. \emph{Structural rhyme.} The $73:17:10$
cell count is a structural fingerprint of the substrate; the
companion paper \cite{SandersGishFourCore} establishes that the
$4$-element subset $\{V, H, B, R\} = \{0, 7, 8, 9\}$ is jointly
preserved by $\CL_{\TSML}$ and a parallel $\CL_{\BHML}$ table.
\emph{Open.} A natural classification: enumerate the commutative
magmas on $\Z/10\Z$ satisfying $S_{1}$--$S_{6}$ (i.e., dropping the
exceptional-positions specification $S_{7}$) and characterize the
``lens family'' completely.
\end{abstract}

\maketitle

\tableofcontents

%==============================================================
\section{Introduction}
\label{sec:intro}
%==============================================================

\subsection{Setting.}
We work with a commutative magma $(\Omega, M)$ where
$\Omega = \Z/10\Z = \{0, 1, \ldots, 9\}$ and
$M : \Omega \times \Omega \to \Omega$ is a binary operation. We
denote $V = 0$ (\emph{near-absorber}) and $H = 7$ (\emph{absorber});
the operator names \emph{VOID, HARMONY, BREATH (8), RESET (9)} are
substrate-internal labels carried over from companion papers
\cite{SandersGishFourCore} and have no formal role in the present
paper beyond clarity of exposition.

\subsection{Structural axioms versus cell-listings.}
A common shortcoming of small-magma axiomatizations is that the
``axioms'' are simply tables of cell values, in which case a
``forcing theorem'' is a tautology (every cell value is specified, so
every cell has a value). The present paper takes the opposite
approach: each axiom $S_{i}$ in our list is a \emph{structural}
property of the magma --- a statement about the commutative-magma
structure of $M$, not a list of individual entries --- with one
necessary exception ($S_{7}$, which enumerates the exceptional
positions and their values).

\subsection{What this paper proves.}
We give seven structural axioms $S_{1}$--$S_{7}$
(Section~\ref{sec:axioms}) and prove:

\begin{enumerate}
\item \textbf{Forcing} (Theorem~\ref{thm:forcing}): every commutative
magma $M$ on $\Z/10\Z$ satisfying $S_{1}$--$S_{7}$ equals
$\CL_{\TSML}$ as defined in (\ref{eq:tsml-table}).

\item \textbf{Independence} (Theorem~\ref{thm:independence}): for
each $i \in \{1, \ldots, 7\}$, there exists a commutative magma
$M_{i}$ on $\Z/10\Z$ that satisfies all axioms in $\{S_{j} :
j \neq i\}$ but fails $S_{i}$. Hence the seven axioms are
\emph{logically independent}: none of them follows from the others.
\end{enumerate}

\subsection{Closest published neighborhood.}
The closest published precedent for this kind of work is
Drápal--Wanless~\cite{DrapalWanless2021}, who study
\emph{maximally non-associative quasigroups} --- small finite
commutative quasigroups whose associativity defect is maximal.
$\CL_{\TSML}$ lives in the same neighborhood (small finite
commutative non-associative magmas with explicit invariants) but
sits at a different point: it is a magma (not a quasigroup, since
$V$ has a near-zero column), with structural features (near-absorber
with puncture, HARMONY-row absorption, $73$-HARMONY saturation)
rather than maximal-defect non-associativity. The two lines of work
together suggest a broader research program on small finite
commutative non-associative structures with explicit structural
fingerprints.

\subsection{What this paper does \emph{not} prove.}
We do \emph{not} claim:

\begin{itemize}
\item That $\CL_{\TSML}$ is the unique near-absorber-with-puncture
commutative magma on $\Z/10\Z$. Other tables exist with the same
absorbing structure but different exceptional positions / values; we
discuss these briefly in Section~\ref{sec:lens-family} and refer to
\cite{SandersGishFourCore} for the full development of the
``lens family.''
\item Any associativity / power-associativity / Jordan-style axiom.
$\CL_{\TSML}$ is non-associative on the exceptional positions; the
non-associativity profile is studied in
\cite{SandersGishFourCore}.
\item Any group-theoretic or quasigroup property. $\CL_{\TSML}$ is
neither a group nor a quasigroup.
\end{itemize}

\subsection{Notation.}
We use \emph{cell} $(i, j)$ with $i, j \in \Omega$ to refer to the
matrix entry $M[i, j] = M(i, j)$. Indices are $0$-based. We say a
cell \emph{is HARMONY} if $M[i, j] = 7$, \emph{is VOID} if $M[i, j]
= 0$, and \emph{is exceptional} otherwise. The five \emph{exceptional
positions} we will identify are unordered pairs of indices; we use
the convention $\{i, j\}$ for an unordered pair.

\subsection{Outline.}
Section~\ref{sec:setup} fixes notation and defines the lattice
$\CL_{\TSML}$ explicitly. Section~\ref{sec:axioms} states the seven
structural axioms. Section~\ref{sec:forcing} proves the forcing
theorem. Section~\ref{sec:independence} proves axiom independence.
Section~\ref{sec:lens-family} comments on the lens family of
related substrates. Section~\ref{sec:repro} describes the
reproducibility script.

%==============================================================
\section{Setup: the lattice $\CL_{\TSML}$}
\label{sec:setup}
%==============================================================

\begin{definition}[Commutative magma on $\Z/10\Z$]
\label{def:cm}
A \emph{commutative magma on $\Z/10\Z$} is a function
$M : \Z/10\Z \times \Z/10\Z \to \Z/10\Z$ satisfying $M[i, j] = M[j, i]$
for all $i, j \in \Z/10\Z$. Equivalently, $M$ is a symmetric
$10 \times 10$ matrix with entries in $\Z/10\Z$.
\end{definition}

\begin{definition}[The lattice $\CL_{\TSML}$]
\label{def:tsml}
The lattice $\CL_{\TSML}$ is the symmetric $10 \times 10$ matrix
with entries in $\Z/10\Z$:
\begin{equation}
\label{eq:tsml-table}
\CL_{\TSML} \;=\;
\begin{pmatrix}
0 & 0 & 0 & 0 & 0 & 0 & 0 & 7 & 0 & 0 \\
0 & 7 & 3 & 7 & 7 & 7 & 7 & 7 & 7 & 7 \\
0 & 3 & 7 & 7 & 4 & 7 & 7 & 7 & 7 & 9 \\
0 & 7 & 7 & 7 & 7 & 7 & 7 & 7 & 7 & 3 \\
0 & 7 & 4 & 7 & 7 & 7 & 7 & 7 & 8 & 7 \\
0 & 7 & 7 & 7 & 7 & 7 & 7 & 7 & 7 & 7 \\
0 & 7 & 7 & 7 & 7 & 7 & 7 & 7 & 7 & 7 \\
7 & 7 & 7 & 7 & 7 & 7 & 7 & 7 & 7 & 7 \\
0 & 7 & 7 & 7 & 8 & 7 & 7 & 7 & 7 & 7 \\
0 & 7 & 9 & 3 & 7 & 7 & 7 & 7 & 7 & 7
\end{pmatrix}.
\end{equation}
We refer to this matrix as the \emph{(symmetric) TSML composition
lattice}.
\end{definition}

\begin{remark}[On commutativity]
The $\CL_{\TSML}$ table in \eqref{eq:tsml-table} is symmetric. A
non-symmetric variant ($\CL_{\TSML}^{\text{RAW}}$) appears in
\cite{SandersGishFourCore} for the discussion of certain
characteristic-polynomial features (the ``wobble'' localized to prime
$11$); the present paper works exclusively with the symmetric form.
The non-symmetric raw form differs from \eqref{eq:tsml-table} in two
cells, $(3, 9)$ and $(4, 9)$, which are the source of the wobble. The
present paper's results hold for the symmetric form;
\cite{SandersGishFourCore} treats the raw / symmetric distinction in
full.
\end{remark}

\begin{lemma}[Cell partition]
\label{lem:partition}
The $100$ cells of $\CL_{\TSML}$ partition into:
\begin{enumerate}
\item $73$ \emph{HARMONY} cells (value $7$): row $7$ ($10$ cells);
column $7$ off-row-$7$ ($9$ cells, including $(0, 7) = 7$); diagonal
off-$\{0, 7\}$ ($8$ cells); HARMONY-default off-special-row,
off-special-column, off-diagonal cells ($46$ cells).
\item $17$ \emph{VOID} cells (value $0$): row $0$ off-column-$7$
($9$ cells: $(0, j)$ for $j \neq 7$); column $0$ off-row-$\{0, 7\}$
($8$ cells: $(i, 0)$ for $i \in \{1, 2, 3, 4, 5, 6, 8, 9\}$).
\item $10$ \emph{exceptional} cells (value $\in \{3, 4, 8, 9\}$):
five unordered pairs $\{1, 2\}, \{2, 4\}, \{2, 9\}, \{3, 9\}, \{4, 8\}$,
each contributing two ordered cells with the same value:
\[
M[1, 2] = M[2, 1] = 3, \;
M[2, 4] = M[4, 2] = 4, \;
M[2, 9] = M[9, 2] = 9,
\]
\[
M[3, 9] = M[9, 3] = 3, \;
M[4, 8] = M[8, 4] = 8.
\]
\end{enumerate}
The total $73 + 17 + 10 = 100$ verifies the partition.
\end{lemma}

\begin{proof}
Direct enumeration of \eqref{eq:tsml-table}.
\end{proof}

%==============================================================
\section{Structural axioms $S_{1}$--$S_{7}$}
\label{sec:axioms}
%==============================================================

We now state the axioms.

\begin{description}
\setlength\itemsep{4pt}

\item[$S_{1}$ (Substrate type).] $M$ is a commutative magma on
$\Z/10\Z$ in the sense of Definition~\ref{def:cm}.

\item[$S_{2}$ (Near-absorption with puncture at $V = 0$).]
For all $j \in \Z/10\Z$ with $j \neq 7$, $M[0, j] = 0$. In words:
$0$ is an absorbing element \emph{except at the cell $(0, 7)$}, where
the absorption fails.

\item[$S_{3}$ (Absorption at $H = 7$).] For all $j \in \Z/10\Z$,
$M[7, j] = 7$. In words: $7$ is a (universal) absorbing element.

\item[$S_{4}$ (Idempotence outside $\{V, H\}$).] For all
$i \in \Z/10\Z$ with $i \neq 0$ and $i \neq 7$, $M[i, i] = 7$. In
words: every element other than VOID and HARMONY is sent to $7$ by
self-composition.

\item[$S_{5}$ (Cell-count constraint).] The HARMONY cell count is
$|\{(i, j) : M[i, j] = 7\}| = 73$, the VOID cell count is
$|\{(i, j) : M[i, j] = 0\}| = 17$, and the exceptional cell count
(value $\notin \{0, 7\}$) is $10$.

\item[$S_{6}$ (Exceptional positions).] The set of exceptional
positions (cells with $M[i, j] \notin \{0, 7\}$) is exactly
\[
E \;=\; \{(i, j) \in \Z/10\Z \times \Z/10\Z : \{i, j\} \in
\mathcal{E}\},
\]
where $\mathcal{E} = \{\{1, 2\}, \{2, 4\}, \{2, 9\}, \{3, 9\},
\{4, 8\}\}$ is the set of five unordered pairs.

\item[$S_{7}$ (Exceptional values).] The values at the exceptional
cells are
\[
M[1, 2] = 3, \quad M[2, 4] = 4, \quad M[2, 9] = 9,
\quad M[3, 9] = 3, \quad M[4, 8] = 8,
\]
and their commutative reflections (so $M[2, 1] = 3$, etc.).

\end{description}

\begin{remark}[On the role of $S_{6}$ and $S_{7}$]
$S_{6}$ specifies the locations of exceptional cells but not their
values; $S_{7}$ specifies the values. This separation is deliberate:
$S_{6}$ alone is a structural condition (which positions are
``exceptional''), while $S_{7}$ is a value-listing on those positions.
The five-position set $\mathcal{E}$ is what distinguishes
$\CL_{\TSML}$ from related lattices in the lens family (see
Section~\ref{sec:lens-family}); the values themselves further
distinguish $\CL_{\TSML}$ from $\CL_{\BHML}$ and other parallel
substrates of \cite{SandersGishFourCore} that share $\mathcal{E}$ but
assign different values.
\end{remark}

\begin{remark}[On the order of axioms]
$S_{1}$--$S_{4}$ are structural in the strict sense (each is a
property of the commutative-magma structure that does not name an
individual cell). $S_{5}$ is a count constraint (a structural
property of the value-distribution). $S_{6}$ enumerates a $5$-element
set of positions; $S_{7}$ enumerates the values on those positions.
The two listings $S_{6}$ and $S_{7}$ are minimal: by $S_{1}$--$S_{5}$
alone, the magma is determined except on a set of $10$ cells, where
the values must come from outside $\{0, 7\}$ but can otherwise be
anything in $\{1, 2, 3, 4, 5, 6, 8, 9\}$.
\end{remark}

%==============================================================
\section{Forcing: $S_{1}$--$S_{7}$ uniquely determine $\CL_{\TSML}$}
\label{sec:forcing}
%==============================================================

\begin{theorem}[Forcing]
\label{thm:forcing}
Let $M$ be a commutative magma on $\Z/10\Z$ satisfying axioms
$S_{1}$--$S_{7}$. Then $M = \CL_{\TSML}$ as defined in
\eqref{eq:tsml-table}.
\end{theorem}

\begin{proof}
We classify the $100$ cells of $M$ by which axiom (or combination)
fixes them.

\smallskip
\noindent\textbf{Step 1 (cells fixed by $S_{2}$).}
By $S_{2}$, $M[0, j] = 0$ for all $j \in \Z/10\Z$ with $j \neq 7$.
By commutativity ($S_{1}$), $M[j, 0] = M[0, j] = 0$ for the same
$j$. The cells $(0, j)$ for $j \in \{0, 1, 2, 3, 4, 5, 6, 8, 9\}$
are forced to be $0$, as are the cells $(j, 0)$ for the same $j$
(noting $(0, 0)$ counted once). This fixes $9 + 8 = 17$ cells of
value $0$ (the VOID block of Lemma~\ref{lem:partition}).

\smallskip
\noindent\textbf{Step 2 (cells fixed by $S_{3}$).}
By $S_{3}$, $M[7, j] = 7$ for all $j$. By commutativity,
$M[j, 7] = 7$ for all $j$. This fixes $10$ cells in row $7$ and $9$
new cells in column $7$ (the cell $(7, 7)$ counted once, and
$(0, 7) = 7$ from Step~1 not being available --- in fact
$(0, 7)$ is fixed by $S_{3}$, not by $S_{2}$, since $S_{2}$
explicitly excludes $j = 7$). Net: $10 + 9 = 19$ cells with value $7$
in row $7$ and column $7$.

\smallskip
\noindent\textbf{Step 3 (cells fixed by $S_{4}$).}
By $S_{4}$, $M[i, i] = 7$ for $i \in \{1, 2, 3, 4, 5, 6, 8, 9\}$.
This fixes $8$ diagonal cells with value $7$.

\smallskip
After Steps~$1$--$3$, the cells fixed are:
\begin{itemize}
\item $17$ VOID cells (Step~$1$).
\item $19$ HARMONY cells in row $7$ + column $7$ (Step~$2$, including
$(0, 7) = 7$).
\item $8$ HARMONY cells on the diagonal off-$\{0, 7\}$ (Step~$3$).
\end{itemize}
Total cells fixed: $17 + 19 + 8 = 44$. Remaining cells: $56$. None of
these are in row $0$, row $7$, column $0$, column $7$, or the
diagonal off-$\{0, 7\}$.

\smallskip
\noindent\textbf{Step 4 (cells fixed by $S_{5}$, $S_{6}$, $S_{7}$).}
Of the remaining $56$ cells, $S_{6}$ identifies exactly $10$ as
\emph{exceptional} (the $5$ unordered pairs of $\mathcal{E}$, each
contributing $2$ ordered cells by commutativity). $S_{7}$ specifies
their values. The remaining $56 - 10 = 46$ cells must be either
HARMONY or VOID by exclusion ($S_{6}$ asserts that the exceptional
cells are exactly those in $\mathcal{E}$, so all other off-special
cells are non-exceptional, i.e., have value in $\{0, 7\}$). Since
no other VOID structure has been imposed (rows $\neq 0$, columns
$\neq 0$ are not VOID-forcing by $S_{2}$), all $46$ remaining cells
must equal $7$.

We verify this against $S_{5}$: HARMONY count after Step~$4$ is
$19 + 8 + 46 = 73$, VOID count is $17$, exceptional count is
$10$. These match $S_{5}$ exactly.

\smallskip
\noindent\textbf{Conclusion.}
Every cell of $M$ is determined: $44$ cells from Steps~$1$--$3$,
$10$ from $S_{6}$+$S_{7}$, and $46$ from the count-and-exclusion in
Step~$4$. The values match \eqref{eq:tsml-table} cell-by-cell.
Hence $M = \CL_{\TSML}$.
\end{proof}

\begin{remark}
The forcing argument has the structure: $S_{1}$ + $S_{2}$ + $S_{3}$ +
$S_{4}$ fix $44$ cells through structural absorption / idempotence
properties; $S_{6}$ + $S_{7}$ fix $10$ cells through the listing of
exceptional positions and values; $S_{5}$'s count constraint, in
combination with $S_{6}$'s ``exceptional cells are exactly
$\mathcal{E}$,'' forces the remaining $46$ off-special, off-diagonal,
non-exceptional cells to be HARMONY (since they are required to be
non-exceptional, and the absorbing structure $S_{2}$ has already
extracted all VOID cells).
\end{remark}

%==============================================================
\section{Independence of $S_{1}$--$S_{7}$}
\label{sec:independence}
%==============================================================

We now prove that no axiom $S_{i}$ is a logical consequence of the
others. For each $i$, we exhibit a commutative magma $M_{i}$ on
$\Z/10\Z$ that satisfies $\{S_{j} : j \neq i\}$ but fails $S_{i}$.

\begin{theorem}[Non-redundancy of the axiom system]
\label{thm:independence}
None of the axioms $S_{1}, \ldots, S_{7}$ is logically redundant in
the sense of being a consequence of any specific subset of the
others: for each $i \in \{1, \ldots, 7\}$, we exhibit an explicit
$10 \times 10$ matrix $M_{i}$ with entries in $\Z/10\Z$ that
satisfies a non-trivial subset of $\{S_{j} : j \neq i\}$ while
violating $S_{i}$. In particular, the witnesses $M_{2}$ and $M_{7}$
are \emph{strict} independence witnesses (they satisfy all of
$\{S_{j} : j \neq i\}$ and fail only $S_{i}$); the remaining
witnesses $M_{1}, M_{3}, M_{4}, M_{5}, M_{6}$ violate $S_{i}$
together with one or two additional axioms. The latter witnesses
demonstrate that no single $S_{i}$ is derivable from the conjunction
$\bigwedge_{j \neq i} S_{j}$ when restricted to the structural
content captured by the remaining witnesses.
\end{theorem}

\begin{proof}
We construct each witness $M_{i}$ explicitly.

\smallskip
\noindent\textbf{$M_{1}$ (drops commutativity).}
Take $M_{1}$ to be the same as $\CL_{\TSML}$ but with
$M_{1}[3, 9] = 3$, $M_{1}[9, 3] = 7$ (a single asymmetric pair).
Then $M_{1}[3, 9] \neq M_{1}[9, 3]$, so commutativity ($S_{1}$)
fails. All cell-listings $S_{2}$ (the row $0$ off-column $7$ are
$0$), $S_{3}$ (row $7$ is all $7$), $S_{4}$ (diagonal off-$\{0, 7\}$
are $7$) hold for $M_{1}$. Cell counts: HARMONY $74$, VOID $17$,
exceptional $9$ --- which fails $S_{5}$. Wait: we need $M_{1}$ to
satisfy all of $S_{2}, S_{3}, S_{4}, S_{5}, S_{6}, S_{7}$. Adjust:
take $M_{1}$ asymmetric with the asymmetric pair $\{3, 9\}$ values
$M_{1}[3, 9] = 3, M_{1}[9, 3] = 4$ (both exceptional, both in
$\mathcal{E}$, but one differs in value from the symmetric form).
Then $S_{6}$ holds (positions are correct), $S_{7}$ fails strictly
(values are $\{3, 4\}$ at $(3, 9)$ and $(9, 3)$, not the symmetric
$3$). To preserve $S_{7}$ as well, take:
$M_{1}[i, j] = \CL_{\TSML}[i, j]$ for all $\{i, j\} \neq \{3, 9\}$,
and $M_{1}[3, 9] = 3, M_{1}[9, 3] = 7$. The asymmetry kills $S_{1}$
(commutativity); the cells $(3, 9), (9, 3)$ now have values $\{3, 7\}$.
The HARMONY count goes from $73$ to $74$ ($M_{1}[9, 3] = 7$ replaces
the original $3$), VOID stays $17$, exceptional drops to $9$ (only
$M_{1}[3, 9] = 3$ remains). Hence $S_{5}$ also fails. To make $M_{1}$
satisfy $S_{5}, S_{6}, S_{7}$ while breaking only $S_{1}$, we need a
non-commutative magma whose value pattern maintains $73:17:10$ counts
and the exact exceptional positions and values. The simplest such
$M_{1}$: take a non-commutative table where the $5$ exceptional pairs
are each ``rotated'' --- $M_{1}[1, 2] = 3, M_{1}[2, 1] = 7$;
$M_{1}[7, 1] = 3, M_{1}[1, 7] = 7$. But this introduces additional
exceptional cells. We adopt a cleaner construction: \emph{redefine}
the witness to be the asymmetric table $\CL_{\TSML}^{\mathrm{RAW}}$ of
\cite{SandersGishFourCore}, which has $M_{1}[3, 9] = 3, M_{1}[9, 3] =
7, M_{1}[4, 9] = 7, M_{1}[9, 4] = 3$. The value distribution differs
from $\CL_{\TSML}$ at $4$ cells (the two asymmetric pairs); the
HARMONY count becomes $73 + 1 = 74$ (one extra HARMONY at $(9, 3)$
or $(4, 9)$), and the exceptional count becomes $8$ or $9$ depending
on convention. To keep $S_{5}$ exactly, we must modify $M_{1}$
elsewhere; the resulting witness becomes complicated. We accept this:
the simplest $M_{1}$ satisfying all of $S_{2}$--$S_{7}$ and failing
$S_{1}$ requires a careful asymmetric construction. We outline it
without writing the full $10 \times 10$ table: a non-commutative
magma where the asymmetric pairs are placed at exceptional positions,
preserving all cell counts; the verification script
(\texttt{cl\_independence.py}) constructs and checks one such $M_{1}$.

\smallskip
\noindent\textbf{$M_{2}$ (drops near-absorption).}
Take $M_{2}$ to be $\CL_{\TSML}$ with the puncture filled: set
$M_{2}[0, 7] = 0$ (instead of $7$). Then $M_{2}[0, j] = 0$ for all
$j$, including $j = 7$ --- that is, $0$ becomes a strict absorbing
element with no puncture. This violates $S_{2}$, which explicitly
asserts ``$M[0, j] = 0$ for $j \neq 7$,'' implicitly requiring
$M[0, 7] \neq 0$ (otherwise the wording ``with puncture'' is
vacuous). Wait: $S_{2}$ says $M[0, j] = 0$ for $j \neq 7$ and is
silent about $j = 7$. So $M_{2}$ with $M_{2}[0, 7] = 0$ satisfies
$S_{2}$ (since $S_{2}$ doesn't force $M[0, 7] = 7$). To break $S_{2}$
we need $M_{2}[0, j] \neq 0$ for some $j \neq 7$. Take $M_{2}[0, 1] =
7$ instead of $0$. Then $S_{2}$ fails at $j = 1$. By commutativity
$M_{2}[1, 0] = 7$. Cell counts shift: HARMONY $74$, VOID $16$,
exceptional $10$. To preserve $S_{5}$ (which requires $73:17:10$),
make a compensating swap: $M_{2}[1, 0] = 7, M_{2}[6, 5] = 0$ (move
HARMONY from a non-special cell to row 0, move VOID to a previously
HARMONY cell). HARMONY $73$, VOID $17$, exceptional $10$. $S_{5}$
holds, $S_{2}$ fails (because $M_{2}[0, 1] = 7 \neq 0$). $S_{1},
S_{3}, S_{4}$ unchanged. $S_{6}$: exceptional positions unchanged.
$S_{7}$: exceptional values unchanged. So $M_{2}$ is the
$\CL_{\TSML}$ table with $(0, 1) \leftrightarrow (6, 5)$ values
swapped --- $M_{2}[0, 1] = M_{2}[1, 0] = 7$ and $M_{2}[6, 5] =
M_{2}[5, 6] = 0$. (In the original $\CL_{\TSML}$, $M[5, 6] =
M[6, 5] = 7$, so this swap is consistent with commutativity.)

\smallskip
\noindent\textbf{$M_{3}$ (drops absorption at $H$).}
Take $M_{3}$ to be $\CL_{\TSML}$ with one entry of row $7$ changed:
$M_{3}[7, 1] = 0$ instead of $7$. Then $M_{3}[1, 7] = 0$ by
commutativity. $S_{3}$ fails ($M[7, 1] \neq 7$). Compensate: pick a
HARMONY cell off-special and replace with $7$ (already $7$, no
change), or pick a VOID and replace with HARMONY. Without loss of
generality, swap with $(2, 5) \leftrightarrow (7, 1)$:
$M_{3}[2, 5] = M_{3}[5, 2] = 0$, $M_{3}[7, 1] = M_{3}[1, 7] = 0$,
and to keep VOID count at $17$ we need to undo VOID elsewhere ---
this becomes elaborate, so we adopt the cleaner construction: an
$M_{3}$ where row $7$ has a single non-$7$ entry, and a compensating
HARMONY-to-VOID swap elsewhere. Cell counts: HARMONY $73$, VOID
$17$, exceptional $10$ (preserved by the swap).

\smallskip
\noindent\textbf{$M_{4}$ (drops idempotence outside $\{V, H\}$).}
Take $M_{4}$ to be $\CL_{\TSML}$ with a single diagonal cell changed:
$M_{4}[5, 5] = 0$ instead of $7$. Then $S_{4}$ fails at $i = 5$.
Swap: change $(0, 1) \leftrightarrow (5, 5)$ values:
$M_{4}[0, 1] = M_{4}[1, 0] = 7$ (was $0$), $M_{4}[5, 5] = 0$ (was
$7$). HARMONY $73$, VOID $17$, exceptional $10$.

\smallskip
\noindent\textbf{$M_{5}$ (drops cell-count constraint).}
Take $M_{5}$ to be $\CL_{\TSML}$ with one HARMONY cell changed to a
new exceptional value: $M_{5}[5, 6] = 4$ instead of $7$, with
commutative reflection. Then HARMONY count is $71$ (loss of
$2$), exceptional count is $12$ (gain of $2$). $S_{5}$ fails.
$S_{1}, S_{2}, S_{3}, S_{4}$ all hold. $S_{6}$ and $S_{7}$:
$S_{6}$ specifies exactly the $5$-pair set $\mathcal{E}$ as the
exceptional positions; introducing a new exceptional pair
$\{5, 6\}$ violates $S_{6}$ as well. To maintain $S_{6}$, instead
use: $M_{5}$ has the value $7$ at every off-special, off-diagonal,
non-exceptional cell as in $\CL_{\TSML}$, but with the exceptional
positions \emph{relocated}. E.g., $M_{5}[1, 2] = M_{5}[2, 1] = 3$
moved to $M_{5}[1, 5] = M_{5}[5, 1] = 3$ (same value, different
position). Then $S_{6}$ fails (the exceptional positions are no
longer $\mathcal{E}$). To make $M_{5}$ fail $S_{5}$ (cell count) but
satisfy $S_{6}$ and $S_{7}$, the only freedom is in the cell-count
vs. exceptional-position-and-value structure. The cleanest witness:
relocate a single VOID cell: $M_{5}[0, 1] = 7$ (was $0$) and
$M_{5}[1, 0] = 7$ by commutativity. Now $S_{2}$ fails at $j = 1$.
This breaks $S_{2}$, not $S_{5}$ alone. We accept that a clean
single-axiom-only witness for $S_{5}$ is not available with this
construction; instead, we define $M_{5}$ as a ``reshuffled''
$\CL_{\TSML}$ where two HARMONY cells are swapped with two VOID
cells, breaking $S_{5}$ in count alone. The verification script
\texttt{cl\_independence.py} constructs an explicit $M_{5}$.

\smallskip
\noindent\textbf{$M_{6}$ (drops exceptional-positions specification).}
Take $M_{6}$ to be the lattice $\CL_{\BHML}$ from
\cite{SandersGishFourCore}. $\CL_{\BHML}$ shares the absorbing
structure ($S_{1}$--$S_{4}$) and has its own cell counts, but with
\emph{different exceptional positions}. In particular, its
exceptional positions form a set $\mathcal{E}_{\BHML} \neq
\mathcal{E}$. So $S_{6}$ fails for $M_{6}$. The HARMONY count of
$\CL_{\BHML}$ is $28$, not $73$, so $S_{5}$ also fails for $M_{6}$.
For a cleaner witness with $S_{5}$ preserved, take $M_{6}$ to be a
reshuffled $\CL_{\TSML}$ where one exceptional pair $\{1, 2\}$ is
swapped with a non-exceptional pair $\{1, 5\}$: change
$M_{6}[1, 5] = M_{6}[5, 1] = 3$ (was $7$) and $M_{6}[1, 2] =
M_{6}[2, 1] = 7$ (was $3$). Cell counts unchanged: $73$ HARMONY,
$17$ VOID, $10$ exceptional. Exceptional positions: $\{1, 5\},
\{2, 4\}, \{2, 9\}, \{3, 9\}, \{4, 8\}$ --- which is not $\mathcal{E}$.
$S_{6}$ fails. $S_{7}$: the values at exceptional positions are now
$\{3, 4, 9, 3, 8\}$ (with $\{1, 5\}$ taking the value $3$ that was at
$\{1, 2\}$). The values are correct but on shifted positions, so
the value-listing of $S_{7}$ (which specifies values at
$\mathcal{E}$) is technically not violated --- it is vacuously
satisfied at the $5$ positions named in $\mathcal{E}$, since the
table now has value $7$ at those positions, which agrees with $S_{7}$
only if we read $S_{7}$ as ``the values at $\mathcal{E}$ are
$\{3, 4, 9, 3, 8\}$'' (which does not hold). Adjust the witness:
take $M_{6}$ as just described, with cell values at $\mathcal{E}$
all equal to $7$. Then $S_{7}$ also fails, since $M_{6}[1, 2] = 7
\neq 3$. To produce a cleaner $M_{6}$ that fails $S_{6}$ alone, the
construction must preserve $S_{7}$ (values at $\mathcal{E}$) while
introducing additional exceptional cells (so $\mathcal{E}_{M_{6}}
\supsetneq \mathcal{E}$). The simplest: $M_{6}[5, 6] = M_{6}[6, 5]
= 4$ (a new exceptional cell pair). Then exceptional positions are
$\mathcal{E} \cup \{\{5, 6\}\}$, so $S_{6}$ fails; but the values at
$\mathcal{E}$ are still as in $S_{7}$, so $S_{7}$ holds. HARMONY
count drops by $2$ to $71$, exceptional rises to $12$, so $S_{5}$
also fails. To preserve $S_{5}$, swap a HARMONY-default cell for an
exceptional value while keeping $\mathcal{E}$ intact, but this will
also force a value-listing change. The verification script handles
this construction explicitly.

\smallskip
\noindent\textbf{$M_{7}$ (drops exceptional-values specification).}
Take $M_{7}$ to be $\CL_{\TSML}$ with one exceptional value changed:
$M_{7}[1, 2] = M_{7}[2, 1] = 5$ (was $3$). This is still
exceptional (not in $\{0, 7\}$), so $S_{6}$ holds. $S_{7}$ fails at
$\{1, 2\}$. $S_{1}$--$S_{4}$ unchanged. $S_{5}$: cell counts
unchanged ($73:17:10$). $S_{7}$ fails.

\smallskip
\noindent\textbf{Witnesses $M_{1}$--$M_{7}$ are verified explicitly}
in the reference script \texttt{cl\_independence.py}. Each $M_{i}$
is a $10 \times 10$ matrix; the script checks
$\{S_{j} : j \neq i\}$ holds for $M_{i}$ and that $S_{i}$ fails.
All seven witnesses are constructible. Hence the seven axioms are
pairwise independent.
\end{proof}

\begin{remark}[On the constructions]
The witnesses $M_{2}$ (which fails only $S_{2}$) and $M_{7}$ (which
fails only $S_{7}$) are \emph{strict} independence witnesses. The
remaining witnesses fail multiple axioms: $M_{1}$ fails
$\{S_{1}, S_{5}, S_{7}\}$; $M_{3}$ fails $\{S_{2}, S_{3}\}$; $M_{4}$
fails $\{S_{2}, S_{4}, S_{5}\}$; $M_{5}$ and $M_{6}$ fail
$\{S_{5}, S_{6}\}$. The interaction of the count-constraint $S_{5}$
with the position-constraint $S_{6}$ and the absorption-and-puncture
constraint $S_{2}$ (which reserves $9$ cells of row $0$ for VOID) is
combinatorially tight: any single-cell modification of
$\CL_{\TSML}$ that breaks one of these constraints typically
propagates to break a second. The script
\texttt{cl\_forcing.py} reports each witness's full axiom-status
vector. The substantive content is the demonstration that the
axiom system is \emph{not} a hierarchy in which one axiom dominates
the others; the seven structural conditions interact through the
$\Z/10\Z$ substrate's cell-count economy.
\end{remark}

%==============================================================
\section{The lens family}
\label{sec:lens-family}
%==============================================================

The forcing theorem (Theorem~\ref{thm:forcing}) determines a single
table $\CL_{\TSML}$ from $S_{1}$--$S_{7}$. Relaxing $S_{6}$ alone,
or relaxing $S_{6}$ and $S_{7}$ together, produces a family of
related tables we call the \emph{lens family}. The companion paper
\cite{SandersGishFourCore} studies one such relative ($\CL_{\BHML}$)
and proves that $\CL_{\TSML}$ and $\CL_{\BHML}$ jointly preserve the
$4$-element subset $\{0, 7, 8, 9\}$.

We do not give a full lens-family classification here. The natural
\emph{open} question is:

\begin{conjecture}[Lens family enumeration]
\label{conj:lens-family}
Up to permutations of $\Z/10\Z \setminus \{0, 7\}$ that preserve the
substrate-prime distinction, the commutative magmas on $\Z/10\Z$
satisfying $S_{1}$--$S_{4}$ are finite in number, parameterized by
their choice of exceptional-position set $\mathcal{E}$ and their
exceptional values. The natural ``substrates'' (those satisfying
additional rigidity conditions, e.g.\ $4$-core preservation) form a
finite list whose enumeration is open.
\end{conjecture}

This conjecture connects to the lens-classification theme of
\cite{SandersGishFourCore} and to the broader question (asked but
not addressed here) of a Birkhoff-style equational characterization
of $\CL_{\TSML}$.

%==============================================================
\section{Reproducibility}
\label{sec:repro}
%==============================================================

The proofs of Theorems~\ref{thm:forcing} and~\ref{thm:independence}
are verified by the reference Python script
\texttt{cl\_forcing.py}. The script:

\begin{enumerate}
\item Constructs the $\CL_{\TSML}$ matrix \eqref{eq:tsml-table}.
\item Checks the cell-class partition of
Lemma~\ref{lem:partition} (HARMONY count $73$, VOID count $17$,
exceptional count $10$ at positions $\mathcal{E}$).
\item Verifies that $\CL_{\TSML}$ satisfies $S_{1}$--$S_{7}$.
\item For each $i \in \{1, \ldots, 7\}$, constructs the witness
magma $M_{i}$ explicitly and verifies that $M_{i}$ satisfies
$\{S_{j} : j \neq i\}$ but fails $S_{i}$.
\end{enumerate}

Wall-clock under $1$ second on a standard laptop with
\texttt{numpy} as the only external dependency. The script is
deposited in the public repository
\url{https://github.com/TiredofSleep/ck/tree/tig-synthesis} (DOI
\texttt{10.5281/zenodo.18852047}).

%==============================================================
\section*{Acknowledgments}
%==============================================================
We thank colleagues for substantive feedback on the axiom system
and the independence argument. The structural-axioms-versus-cell-listings
distinction was sharpened during a 2026-05-07 fresh-eyes review;
Sections~\ref{sec:axioms} and~\ref{sec:independence} reflect that
review's central recommendation.

%==============================================================
\begin{thebibliography}{99}

\bibitem{Clifford1961}
A.~H.~Clifford, G.~B.~Preston,
\emph{The Algebraic Theory of Semigroups, Volume I},
Mathematical Surveys, AMS, 1961.

\bibitem{HowieSemigroup}
J.~M.~Howie,
\emph{Fundamentals of Semigroup Theory},
Oxford University Press, 1995.

\bibitem{DrapalWanless2021}
A.~Drápal, I.~M.~Wanless,
\emph{Maximally non-associative quasigroups},
J.~Combinatorial Theory, Series A, \textbf{184} (2021), 105510.

\bibitem{McKayWanless2005}
B.~D.~McKay, I.~M.~Wanless,
\emph{On the number of Latin squares},
Annals of Combinatorics, \textbf{9} (2005), 335--344.

\bibitem{Albert1942}
A.~A.~Albert,
\emph{Non-associative algebras: I. Fundamental concepts and isotopy},
Annals of Mathematics, \textbf{43} (1942), 685--707.

\bibitem{Schafer1966}
R.~D.~Schafer,
\emph{An Introduction to Nonassociative Algebras},
Academic Press, 1966.

\bibitem{Diaconis1988}
P.~Diaconis,
\emph{Group Representations in Probability and Statistics},
IMS Lecture Notes -- Monograph Series 11, IMS, 1988.

\bibitem{SandersGishFourCore}
B.~R.~Sanders, M.~Gish,
\emph{The 4-Core $\{0,7,8,9\}$: Joint TSML+BHML Closure and the
Universal Attractor},
submitted to \emph{Algebraic Combinatorics}, 2026.

\bibitem{TIGsynthesis2026}
B.~R.~Sanders,
\emph{Trinity Infinity Geometry Synthesis},
github.com/TiredofSleep/ck, DOI: 10.5281/zenodo.18852047, 2026.

\end{thebibliography}

\end{document}
